\scanpage{49}

\margin{de variétés isogènes à $E\times F$ ! ]}

Plus précisément, \unsure{tte courbe} isogène à $E\times F$ \deleted{\unsure{est}} --- sur un corps alg. clos --- est \deleted{isogène} \inserted{\uline{isomorphe}} à un produit $E'\times F'$. Donc les variétés isogènes à un produit font une sous-famille loc fermée de $\unsure{E|N}$, image de l'ouvert $M'_1$ de la variété modulaire $M_1$ des courbes elliptiques, formé de celles qui sont ordinaires, par les applications $E\mapsto E\times F$, où $F$ parcourt le \unsure{nb} fini de types de courbes ell. hypersingulières. \inserted{D'où un ouvert non vide $U$ de $N$ \deleted{\unsure{simple}} formé de variétés \uline{simples}.}

\deleted{Notons que si $A$ est une variété abélienne $\unsure{\in N}$ \unsure{qui est radicalement} \uline{simple} (i.e. qui n'est pas \unsure{le produit de polarisées} de dim $1$) elle n'est pas isogène à \inserted{un produit} $E\times F$}

D'autre part, si $A$ est \uline{simple} de la structure formelle envisagée, alors on sait comme dessus que \deleted{\unsure{les}} les $A'$ isogènes à $A$ forment une famille \uline{discrète} (bien qu'infinie), de façon précise si $A$ est définissable sur $k$, alors $A'$ est définissable sur une extension \uline{finie} de $k$. Si nous prenons donc pts de \deleted{$\unsure{N}$} $U\subset N$ \deleted{\unsure{l'un absolu et générique, l'on}} ayant des degrés de transcendance distincts sur le corps premier $\mathbb F_p$ (les degrés possibles sont $0,1,2$) on trouve des variétés \uline{simples} \uline{non isogènes}, de $\overline{\mathrm m}$ grp formel $\overline G_m\times G_{1,1}$.

\margin{59}
